\section{Basic Concepts of Algorithms}
\label{sec:algorithms}

Our main matrix combines four classifiers with four imbalance-handling
strategies, chosen to span complementary families of decision-boundary
assumptions.

\subsection{Classifiers}
\textbf{Decision Tree (DT).} We grow a CART tree with the entropy
criterion and cost-complexity pruning at $\alpha=0.01$
\cite{breiman1984classification}; the scheme differs from Weka's J48
implementation of C4.5 \cite{quinlan1993c45,frank2016weka}, and the
light pruning keeps the tree interpretable while limiting overfitting.

\textbf{Random Forest (RF).} The forest averages 100 de-correlated
trees grown on bootstrap samples \cite{breiman2001random}; in our
plain baselines it handles imbalance best on sick and stays
competitive with LR on credit-g.

\textbf{Logistic Regression (LR).} A linear decision boundary passed
through a sigmoid, trained by regularised log-loss; the boundary is
the most sensitive to how nominal attributes are encoded.

\textbf{Support Vector Classifier (SVC).} A maximum-margin classifier
with an RBF kernel \cite{cortes1995support}. SVC has no native
probabilities, so \texttt{probability=True} wraps an internal Platt
calibration \cite{platt1999probabilistic} fitted by scikit-learn's own
internal cross-validation, a loop separate from both the 5-fold
threshold search of Section~\ref{ssec:strategies} and the 10-fold
evaluation. We route all four models through \texttt{predict\_proba}
so that AUC and threshold moving share one uniform interface at the
cost of extra training time.

\subsection{Imbalance-Handling Strategies}
\label{ssec:strategies}
\textbf{plain} keeps the empirical class distribution and each
classifier's default \texttt{predict} rule, and serves as the
reference arm for every comparison. For the RBF SVC this rule is the
sign of the decision function rather than a probability threshold.

\textbf{SMOTE} over-samples the minority class by interpolating
between minority neighbours and their nearest neighbours
\cite{chawla2002smote}; the sampler sits inside the \texttt{imblearn}
pipeline, so synthesis is confined to the training folds. Because
synthesis happens in the encoded feature space, interpolated one-hot
or ordinal coordinates need not correspond to a natural category
value, a limitation of the frozen design.

\textbf{Cost-sensitive learning} reweights the training loss
\cite{elkan2001foundations}: on credit-g we apply the official 5:1
cost matrix, $5\times$FN against $1\times$FP, via
\texttt{class\_weight}; sick has no official costs, so an
equal-weight control is used there.

\textbf{Threshold moving} keeps training identical to plain and moves
only the operating point: a threshold $\tau$ is selected on inner
5-fold out-of-fold probabilities of the training fold over the grid
$0.05$ to $0.95$, maximising F1, then locked. It is the only strategy
that never touches training.

\subsection{Ensemble Methods: Bagging and Boosting}
\label{ssec:ensemble-concepts}
Beyond the four base classifiers we conduct an exploratory ensemble
study. \textbf{Bagging} aggregates the majority votes of base
learners grown on bootstrap samples \cite{breiman1996bagging}, and a
random forest adds random feature subspaces to the same idea.
\textbf{Boosting}, here AdaBoost, builds learners sequentially and
re-weights the instances the previous ones misclassified
\cite{freund1997decision}, an implicit form of cost-sensitive learning
on hard, often minority, examples. Both ensembles use the same pruned
decision tree as the base learner, with $n_{\text{estimators}}=100$
matching the random-forest budget for comparability.
